% Options for packages loaded elsewhere
\PassOptionsToPackage{unicode}{hyperref}
\PassOptionsToPackage{hyphens}{url}
\PassOptionsToPackage{dvipsnames,svgnames,x11names}{xcolor}
%
\documentclass[
  11pt,
  a4paper,
]{ltjsarticle}
\usepackage{amsmath,amssymb}
\usepackage{iftex}
\ifPDFTeX
  \usepackage[T1]{fontenc}
  \usepackage[utf8]{inputenc}
  \usepackage{textcomp} % provide euro and other symbols
\else % if luatex or xetex
  \usepackage{unicode-math} % this also loads fontspec
  \defaultfontfeatures{Scale=MatchLowercase}
  \defaultfontfeatures[\rmfamily]{Ligatures=TeX,Scale=1}
\fi
\usepackage{lmodern}
\ifPDFTeX\else
  % xetex/luatex font selection
\fi
% Use upquote if available, for straight quotes in verbatim environments
\IfFileExists{upquote.sty}{\usepackage{upquote}}{}
\IfFileExists{microtype.sty}{% use microtype if available
  \usepackage[]{microtype}
  \UseMicrotypeSet[protrusion]{basicmath} % disable protrusion for tt fonts
}{}
\makeatletter
\@ifundefined{KOMAClassName}{% if non-KOMA class
  \IfFileExists{parskip.sty}{%
    \usepackage{parskip}
  }{% else
    \setlength{\parindent}{0pt}
    \setlength{\parskip}{6pt plus 2pt minus 1pt}}
}{% if KOMA class
  \KOMAoptions{parskip=half}}
\makeatother
\usepackage{xcolor}
\usepackage[margin=24mm]{geometry}
\setlength{\emergencystretch}{3em} % prevent overfull lines
\providecommand{\tightlist}{%
  \setlength{\itemsep}{0pt}\setlength{\parskip}{0pt}}
\setcounter{secnumdepth}{5}
\ifLuaTeX
  \usepackage{selnolig}  % disable illegal ligatures
\fi
\IfFileExists{bookmark.sty}{\usepackage{bookmark}}{\usepackage{hyperref}}
\IfFileExists{xurl.sty}{\usepackage{xurl}}{} % add URL line breaks if available
\urlstyle{same}
\hypersetup{
  colorlinks=true,
  linkcolor={blue},
  filecolor={Maroon},
  citecolor={Blue},
  urlcolor={blue},
  pdfcreator={LaTeX via pandoc}}

\author{}
\date{}

\begin{document}

\hypertarget{ux89e3ux7b54-07-ux884cux5217ux5206ux89e3}{%
\section{解答 07 ---
行列分解}\label{ux89e3ux7b54-07-ux884cux5217ux5206ux89e3}}

\hypertarget{section}{%
\subsection{1}\label{section}}

第2行から第1行の2倍を引くので

\[
L=\begin{bmatrix}1&0\\2&1\end{bmatrix},\quad
U=\begin{bmatrix}2&1\\0&1\end{bmatrix}.
\]

\hypertarget{section-1}{%
\subsection{2}\label{section-1}}

第一列 \texttt{a\_1=(1,1,0)\^{}T} より \texttt{q\_1=a\_1/√2}。第二列から
\texttt{q\_1} 方向を引いて正規化すれば \texttt{q\_2}
が得られます。\texttt{R=Q\^{}TA} として係数を求められます。ここでは、QR
が「列空間の正規直交基底 Q」と「その基底での列座標
R」を分けることが重要です。

\hypertarget{section-2}{%
\subsection{3}\label{section-2}}

\[
A^TA=(UΣV^T)^T(UΣV^T)=VΣ^TU^TUΣV^T.
\]

\texttt{U\^{}TU=I} より結論を得ます。

\hypertarget{section-3}{%
\subsection{4}\label{section-3}}

\texttt{V\^{}T} と \texttt{U}
は可逆な直交変換なのでランクを変えません。したがって
\texttt{rank(A)=rank(Σ)}。対角型の \texttt{Σ}
のランクは非零対角成分、つまり非零特異値の個数です。

\hypertarget{section-4}{%
\subsection{5}\label{section-4}}

\texttt{v\_i} は入力楕円体の主軸になる入力方向、\texttt{σ\_i}
はその方向の伸縮率、\texttt{u\_i} は伸縮後の出力方向です。式では
\texttt{Av\_i=σ\_iu\_i} です。

\hypertarget{section-5}{%
\subsection{6}\label{section-5}}

固有値分解は主に正方行列に対し、同じ固有ベクトル基底で入力・出力を表し、固有値は負・複素にもなり得ます。SVD
は任意の長方形行列に存在し、入力 \texttt{V} と出力 \texttt{U}
に別の直交基底を使い、特異値は常に非負です。

\end{document}
